% Lean: EconLean/Macro/Fiscal/ChamleyJudd.lean
\chapter{Optimal Capital Taxation}
\label{chap:fiscal}

Sources: Chamley (1986); Judd (1985); Straub--Werning (2020); Ljungqvist--Sargent ch.~16.

\begin{definition}[Ramsey first-order conditions]
\label{def:cj_foc}
\lean{Econ.ChamleyJudd.IsRamseyFOC}
\leanok
$uc_t = \beta\,uc_{t+1}(1 + (1-\tau_{t+1})r_{t+1})$ (household Euler) and $Wc_t = \beta\,Wc_{t+1}(1+r_{t+1})$
(planner's capital first-order condition on the pseudo-utility).
\end{definition}

\begin{theorem}[Chamley--Judd]
\label{thm:cj}
\lean{Econ.ChamleyJudd.chamley_judd}
\leanok
\uses{def:cj_foc}
If $uc, Wc, r, \tau$ converge with $uc^*, Wc^* > 0$ and $r^* \ne 0$, then $\tau^* = 0$.
\end{theorem}
\begin{proof}
\leanok
Pass to the limit in both equations: $\beta(1+(1-\tau^*)r^*) = 1 = \beta(1+r^*)$, so $\beta\tau^*r^* = 0$.
\end{proof}

\begin{theorem}[Straub--Werning caveat]
\label{thm:sw}
\lean{Econ.ChamleyJudd.ratio_unbounded, Econ.ChamleyJudd.wedge_unbounded}
\leanok
\uses{def:cj_foc}
If $\underline\tau \le \tau_t \le 1$ with $\underline\tau>0$ and $r_t \ge \underline r > 0$, then $uc_t/Wc_t$ is
unbounded, so the convergence hypothesis of Theorem~\ref{thm:cj} fails. A constant tax $\tau>0$ is an implicit tax
$((1+r)/(1+(1-\tau)r))^T$ on date-$T$ consumption, which is unbounded.
\end{theorem}
\begin{proof}
\leanok
Dividing the two conditions, the ratio grows each period by $(1+r_{t+1})/(1+(1-\tau_{t+1})r_{t+1})$, which is at least
$q = (1+\underline r)/(1+(1-\underline\tau)\underline r) > 1$. Hence the ratio is at least $q^t$ times its initial
value.
\end{proof}

\section{The Ramsey problem in primal form}

\begin{definition}[Primal Ramsey problem]
\label{def:ramsey_primal}
\lean{Econ.RamseyPrimal.Problem, Econ.RamseyPrimal.Problem.Feasible, Econ.RamseyPrimal.Problem.value, Econ.RamseyPrimal.Problem.IsOptimal, Econ.RamseyPrimal.Problem.varK, Econ.RamseyPrimal.Problem.varC}
\leanok
For a given multiplier on implementability, maximize $\sum_t\beta^tW(c_t)$ subject to $k_{t+1} = F(k_t) - c_t - g_t$ from a
given $k_0$, over paths with $c, k \ge 0$ and summable value.
\end{definition}

\begin{theorem}[Planner's Euler equation and Chamley--Judd]
\label{thm:ramsey_primal}
\lean{Econ.RamseyPrimal.Problem.feasible_var, Econ.RamseyPrimal.Problem.varC_eq, Econ.RamseyPrimal.Problem.value_var, Econ.RamseyPrimal.Problem.planner_euler, Econ.RamseyPrimal.Problem.isRamseyFOC, Econ.RamseyPrimal.Problem.chamley_judd_primal, Econ.RamseyPrimal.Problem.nonvacuous}
\leanok
\uses{def:ramsey_primal, def:cj_foc, thm:cj}
At an optimum with $c_t, c_{t+1}, k_{t+1} > 0$, $W'(c_t) = \beta W'(c_{t+1})F'(k_{t+1})$. If the allocation is decentralized by capital taxes, the
conditions of Definition~\ref{def:cj_foc} hold, so convergence implies a zero limiting capital tax.
\end{theorem}
\begin{proof}
\leanok
Save $\varepsilon$ more at $t$ and consume the proceeds at $t+1$. For small $|\varepsilon|$ this is feasible
(interiority and continuity of $F$), and it changes the value only at two
dates, so $J(\varepsilon) \le J(0) = 0$. The derivative at the interior maximum vanishes:
$-\beta^tW'(c_t) + \beta^{t+1}W'(c_{t+1})F'(k_{t+1}) = 0$. For non-vacuity, with $\beta = 1/2$, $W = \mathrm{id}$,
$F(k) = 2k$, the plan $c = k = 1$ is optimal: telescoping gives $\sum_t2^{-t}c'_t \le 2k'_0$.
\end{proof}

\section{Ricardian equivalence and tax smoothing}

\begin{theorem}[Ricardian equivalence]
\label{thm:ricardian}
\lean{Econ.Ricardian.pv, Econ.Ricardian.budget, Econ.Ricardian.GovBudget, Econ.Ricardian.pv_sub, Econ.Ricardian.budget_eq, Econ.Ricardian.ricardian, Econ.Ricardian.budgetNB, Econ.Ricardian.not_ricardian_of_constraint}
\leanok
Two tax paths that both satisfy the government's budget constraint give the same household budget set, and hence
the same optimal consumption. With a no-borrowing constraint, a debt-financed tax cut can change optimal
consumption.
\end{theorem}
\begin{proof}
\leanok
The budget depends on taxes only through $\mathrm{PV}(\tau) = \mathrm{PV}(g)$. For the counterexample, take income
$(0,2)$, taxes $(0,0)$ versus $(-1,1)$, and $U = -(c_0-1)^2 - (c_1-1)^2$: after the cut $(1,1)$ gives $U = 0$,
while without it $c_0 \le 0$ forces $U \le -1$.
\end{proof}

\begin{theorem}[Tax smoothing]
\label{thm:tax_smoothing}
\lean{Econ.TaxSmoothing.cost, Econ.TaxSmoothing.Feasible, Econ.TaxSmoothing.ann, Econ.TaxSmoothing.ann_pos, Econ.TaxSmoothing.smooth, Econ.TaxSmoothing.smooth_feasible, Econ.TaxSmoothing.cost_decomp, Econ.TaxSmoothing.smooth_optimal, Econ.TaxSmoothing.smooth_unique}
\leanok
With $\beta R = 1$, the constant tax $B/\sum_t\beta^t$ uniquely minimizes $\sum_t\beta^t\tau_t^2$ subject to
$\sum_t\beta^t\tau_t = B$.
\end{theorem}
\begin{proof}
\leanok
Cost equals the smooth cost plus the cost of the deviation, because the cross term is proportional to
$\sum_t\beta^t(\tau_t - \bar\tau) = 0$.
\end{proof}
